\section{Prescribed factors and prime-support transport}
\label{sec:interfaces}

After the initial Type-I reduction write a contributing modulus as
$d=qr$, where $q\asymp Q$, $r\asymp R$, and put

\begin{equation}\label{eq:H}
 H=\frac{x^\eta RQ^2}{q_0M}.
\end{equation}

The witness $u\mid d/r$, after its common part with $q_0$ is removed,
gives $q_1=u_1v_1$ and the ranges

\begin{align}
 q_0^{-2}x^{-\delta-5\eta}\frac QH\ll U
 &\ll q_0^{-1}x^{-5\eta}\frac QH,\notag\\
 x^{5\eta}H\ll V&\ll q_0x^{\delta+5\eta}H,
 \qquad UV\asymp Q/q_0.
\label{eq:UV}
\end{align}

The second witness and $d\asymp QR$ give

\[
 d_1\asymp\frac{x^{2-\gamma-52\eta}}{Q^4R^2}
 \asymp\frac{N}{q_0^2x^{50\eta}H^2}.
\]

After the standard subdivision into intervals shorter by $x^{5\eta}$,

\begin{equation}\label{eq:Delta1}
 \frac{N}{q_0^2x^{\delta+55\eta}H^2}
 \ll\Delta_1\ll
 \frac{N}{q_0^2x^{55\eta}H^2}.
\end{equation}

\subsection{Reversible small-prime preprocessing}

Put

\[
 D_0=\exp((\log x)^{1/3}),\qquad
 S_x=\exp((\log x)^{2/3}).
\]

\begin{lemma}[Preprocessing]\label{lem:preprocessing}
Let $d$ be squarefree and let $r_0,u_0,d_{1,0}$ witness membership in
$\DS(x)$.  Define

\[
 s(d)=\prod_{\substack{p\mid d\\p\le D_0}}p,\qquad
 q_{\rm raw}=d/r_0,\qquad t=(q_{\rm raw},s(d)).
\]

If $s(d)\le S_x$, then

\[
 q=q_{\rm raw}/t,\qquad r=r_0t
\]

gives $d=qr$ and $(q,P(D_0))=1$.  There exist $u'\mid q$ and
$d_{1,0}\mid r$ satisfying, after dyadic localization,

\begin{align}
 S_x^{-1}\frac{x^{1-\gamma-6\eta-\delta}}{QR}\ll u'
 &\ll\frac{x^{1-\gamma-6\eta}}{QR},\label{eq:pre-u}\\
 S_x^{-2}\frac{x^{2-\gamma-52\eta-\delta}}{Q^4R^2}\ll d_{1,0}
 &\ll\frac{x^{2-\gamma-52\eta}}{Q^4R^2}.
\label{eq:pre-d1}
\end{align}
\end{lemma}

\begin{proof}
Squarefreeness makes the displayed factorization exact.  Every small prime
formerly in $q_{\rm raw}$ divides $t$, so the new $q$ is free of primes at
most $D_0$.  Set $u'=u_0/(u_0,t)$.  Division by at most $S_x$ proves
\cref{eq:pre-u}.  Since $r_0=r/t$, insertion into
\cref{eq:S52-d1} produces exactly the factor $t^{-2}$ and proves
\cref{eq:pre-d1}.
\end{proof}

The exceptional set $s(d)>S_x$ is negligible by the standard
small-prime-tail estimate in the alternative Type-I reduction.  Restoring
the preprocessing changes \cref{eq:UV,eq:Delta1} only by
$S_x^{O(1)}=x^{o(1)}$.  Every later expression has fixed algebraic depth,
so this loss never iterates with $x$.

Define the independent families

\begin{equation}\label{eq:families}
\begin{aligned}
\Qfam(Q,R)=\{q\in[Q,2Q]:{}&(q,P(D_0))=1,\\[-1mm]
&q\text{ has a divisor in \cref{eq:pre-u}}\},\\
\Rfam(Q,R)=\{r\in[R,2R]:{}&
r\text{ has a divisor in \cref{eq:pre-d1}}\}.
\end{aligned}
\end{equation}

Each nonexceptional modulus lies in one selected pair from
$\Qfam\times\Rfam$.  A witness is selected before any Cartesian
enlargement or Cauchy copy is made.

\begin{lemma}[First-factor extraction]\label{lem:first-factor}
Let $q^\sharp=q_0q_1$ be squarefree and let $u\mid q^\sharp$.  If

\[
 a_0=(u,q_0),\qquad u_1=u/a_0,\qquad v_1=q_1/u_1,
\]

then $q_1=u_1v_1$ and $u/q_0\le u_1\le u$.  Consequently the ranges
\cref{eq:UV} hold.
\end{lemma}

\begin{proof}
Prime supports in a squarefree number are sets.  Removing from $u$ the
support shared with $q_0$ leaves a divisor of $q_1$; $v_1$ is its unique
complement.  The inequalities follow from $a_0\le q_0$.
\end{proof}

\begin{lemma}[Second-factor extraction]\label{lem:second-factor}
Choose one witness $d_1\mid r$, write $r=d_1r_1$, and localize
$d_1\asymp\Delta$.  Then

\[
 S_x^{-2}\frac{N}{q_0^2x^{\delta+50\eta}H^2}\ll\Delta
 \ll\frac{N}{q_0^2x^{50\eta}H^2}.
\]

An $O(x^{5\eta})$ smooth cover at scale
$\Delta_1=\Delta/x^{5\eta}$ gives \cref{eq:Delta1}; moreover $d_1,r_1$
are squarefree and coprime.
\end{lemma}

\begin{lemma}[Prime-support transport]\label{lem:prime-support}
Suppose $q_0u_1v_ir$ and $q_0q_2r$ are squarefree for $i=1,2$, and
$(u_1v_1v_2,q_0q_2)=1$.  Let

\[
 r=d_1r_1,\qquad g=(v_1,v_2),\qquad v_i=gv_i^*.
\]

Then the supports of

\[
 d_1,r_1,q_0,u_1,g,v_1^*,v_2^*,q_2
\]

are pairwise disjoint.  Hence

\begin{equation}\label{eq:m}
 m=r_1q_0u_1[v_1,v_2]q_2
 =r_1q_0u_1gv_1^*v_2^*q_2
\end{equation}

is squarefree, $(d_1,m)=1$, and $q_0q_3$ is squarefree for every
$q_3\mid m/q_0$.
\end{lemma}

\begin{proof}
For squarefree integers, divisibility is inclusion of prime supports, gcd is
intersection, lcm is union, and coprimality is disjointness.  The hypotheses
separate all listed supports except the permitted overlap of $v_1,v_2$;
extracting $g$ partitions that last overlap.  Splitting $r=d_1r_1$ proves
the remaining assertions.
\end{proof}

\begin{lemma}[Valuation saturation]\label{lem:valuation}
Let

\[
 m^*=m^{\lfloor2\log x\rfloor},\quad
 w_2=(\lambda,m^*)=(\widetilde\lambda,m^*),\quad
 \lambda^*=(\lambda,\widetilde\lambda).
\]

If $|\lambda|,|\widetilde\lambda|<x$, then, for large $x$,

\[
 (\lambda/w_2,m)=(\widetilde\lambda/w_2,m)
 =(\lambda/\lambda^*,m)=(\widetilde\lambda/\lambda^*,m)=1.
\]
\end{lemma}

\begin{proof}
For every $p\mid m$, one has
$p^{\lfloor2\log x\rfloor}>x$.  Thus the gcd with $m^*$ captures the full
$p$-adic part of each integer.  Removing $w_2$, or the larger common factor
$\lambda^*$, leaves a unit modulo $m$.
\end{proof}

\begin{lemma}[Gcd average without a range hypothesis]
\label{lem:gcd-average}
For $m\ge1$, $K,T>0$, and integers $A\ne0,B$,

\begin{equation}\label{eq:gcd-average}
 \sum_{\substack{|k|\le K\\(Ak+B,m)\le T}}(Ak+B,m)
 \le\tau(m)\bigl(2(A,m)K+T\bigr).
\end{equation}
\end{lemma}

\begin{proof}
Use $(u,m)\le\sum_{c\mid(u,m)}c$, interchange the two finite sums, and
enlarge equality of the gcd to $c\mid Ak+B$.  If soluble, this congruence
is one class modulo $c/(A,c)$ and contributes at most
$2K(A,c)/c+1$.  Multiplication by $c$ gives at most
$2(A,m)K+T$ for every $c\le T$.  There are at most $\tau(m)$ such divisors.
\end{proof}
